\chapter{Your first cut}
\label{ch:firstcut}

This is the chapter to do with your hands on the machine. We will cut and engrave
a small plywood tag: one cut operation for the outline and the hanging hole, one
engrave operation for a line of lettering. Nothing here is specific to a
particular machine except the numbers, and the numbers are the part you will
adjust.

\begin{figure}[htbp]
\centering
\begin{tikzpicture}[scale=0.95]
\draw[line width=0.9pt,mkred] (0,0) rectangle (6,3);
\draw[line width=0.9pt,mkred] (1.6,1.5) circle (0.28);
\draw[mkblue,line width=0.6pt] (2.5,1.7) -- (5.4,1.7);
\draw[mkblue,line width=0.6pt] (2.5,1.2) -- (5.4,1.2);
\draw[mkblue,line width=0.6pt] (2.5,0.7) -- (4.6,0.7);
\node[font=\sffamily\scriptsize,color=mkgrey] at (1.6,0.85) {hole};
\node[font=\sffamily\scriptsize,color=mkgrey] at (3.9,-0.35) {cut: outline + hole (red)};
\node[font=\sffamily\scriptsize,color=mkblue] at (3.95,2.15) {engrave: text (blue)};
\draw[{Stealth[length=2mm]}-{Stealth[length=2mm]},mkdark] (0,-0.75) -- (6,-0.75);
\node[font=\sffamily\scriptsize,color=mkdark,fill=white] at (3,-0.75) {60 mm};
\draw[{Stealth[length=2mm]}-{Stealth[length=2mm]},mkdark] (6.7,0) -- (6.7,3);
\node[font=\sffamily\scriptsize,color=mkdark,fill=white,rotate=90] at (6.7,1.5) {30 mm};
\end{tikzpicture}
\caption{The finished job: a 60\,$\times$\,30\unit{mm} tag. Two operations share
one drawing --- a red \mkseries{Cut} operation for the outline and the 4\unit{mm}
hole, and a blue \mkseries{Engrave} operation for the text. Colour is how you
keep track of which is which.}
\label{fig:first-tag}
\end{figure}

\section{Before you press Start}

\begin{machinebox}{The five-minute setup}
\begin{enumerate}[label=\arabic*.,leftmargin=1.4em]
  \item \textbf{Focus.} Set the lens height for your material with the machine's
        own gauge or jig. Software cannot fix focus.
  \item \textbf{Air assist on} (if fitted). It clears smoke from the beam path and
        reduces flare.
  \item \textbf{Secure the material.} Tape it, magnet it, or clamp it. A workpiece
        that shifts mid-job produces scrap and possibly a fire.
  \item \textbf{Focus test on scrap.} If you have never cut this material, treat
        your first job as the test. Keep a fire extinguisher within reach and
        never leave the machine running unattended.
  \item \textbf{Know your escape.} Locate the pause and stop controls --- on the
        machine, and in the \panel{Laser-Control} pane (\btn{Pause},
        \btn{Stop}).
\end{enumerate}
\end{machinebox}

\section{Step 1 --- Start a project and check the bed}

\begin{walkthrough}{A new project}
\begin{enumerate}[label=\arabic*.,leftmargin=1.4em]
  \item Start MeerK40t and make sure your machine is the active device (the red
        row in the \panel{Devices} pane, Chapter~\ref{ch:connect}).
  \item \menu{File>New} (\key{Ctrl+N}) clears operations, elements and notes. You
        now have an empty document.
  \item Press \key{Ctrl+B} (zoom to bed). The white rectangle you see is your
        machine's working area, drawn from the device's \emph{Width} and
        \emph{Height} settings. If the proportions look wrong, stop and fix the
        device configuration now --- everything downstream depends on it.
\end{enumerate}
\end{walkthrough}

\section{Step 2 --- Draw the shape}

You can draw the tag here or import it; both are normal. Drawing it makes the
later steps easier to see.

\begin{walkthrough}{Draw a rectangle, a hole and some text}
\begin{enumerate}[label=\arabic*.,leftmargin=1.4em]
  \item Click \btn{Rectangle} in the \panel{Tools} icon bar. Click once near the
        top-left of the bed for the first corner, and again for the opposite
        corner. A rectangle appears and is selected.
  \item Open the \panel{Position} pane (\menu{Panes>Editing>Position}). Set
        \btn{Width} to \cmd{60mm} and \btn{Height} to \cmd{30mm}, pressing
        \key{Enter} in each field. Set \btn{X} and \btn{Y} to place it, for
        example \cmd{20mm} and \cmd{20mm}.
  \item Click \btn{Circle} in the \panel{Tools} bar and click near the left end of
        the rectangle to place a circle. In the \panel{Position} pane set both
        \btn{Width} and \btn{Height} to \cmd{4mm} (a circle with equal width and
        height), then set \btn{X} to \cmd{26mm} and \btn{Y} to \cmd{33mm} so it
        sits 8\unit{mm} from the end and centred vertically.
  \item Click \btn{Text} in the \panel{Tools} bar and click to the right of the
        hole. With the new text element selected, open \menu{Tools>Properties}
        and type your lettering, for example the word \cmd{WORKSHOP}. Set the
        height so it fits in the remaining width.
\end{enumerate}
\end{walkthrough}

\noindent If you prefer to type rather than click, the console does the same work.
\cmd{rect 20mm 20mm 60mm 30mm} creates the rectangle, \cmd{circle 26mm 33mm 4mm}
the hole, and \cmd{text WORKSHOP} a text element. Console-created elements are
selected automatically, so \cmd{scale}, \cmd{position} and friends apply to them
immediately --- see Chapter~\ref{ch:console}.

\begin{tipbox}[Get the size right before the burn, not after]
Importing from Inkscape or Illustrator usually brings a scale with it: the
drawing arrives at whatever size the file says, which is rarely the size you
want. Press \key{Ctrl+B}, look at the bed, and use the \panel{Position} pane to
set the real width. The dimension lines drawn on the scene update as you type.
\end{tipbox}

\section{Step 3 --- Create the two operations}

Operations are recipes, and this job needs two: one to cut, one to engrave.

\begin{walkthrough}{Cut and engrave}
\begin{enumerate}[label=\arabic*.,leftmargin=1.4em]
  \item Right-click the \emph{Operations} branch in the \panel{Tree} and choose
        \menu{Append operation>Append Cut}. A red operation appears. Rename it if
        you like --- a name such as \emph{Cut 3mm ply} is worth more than
        \emph{Cut} when you have six operations.
  \item Right-click \emph{Operations} again and choose \menu{Append
        operation>Append Engrave}. A blue operation appears below it.
  \item The order of the list is the burning order. Cutting last is usually what
        you want: engrave the text first, while the part is still held by the
        surrounding material. Drag one operation above the other if you need to
        reorder; dragging in the tree changes burn order.
\end{enumerate}
\end{walkthrough}

\section{Step 4 --- Assign the artwork}

\begin{walkthrough}{Put each shape in the right operation}
\begin{enumerate}[label=\arabic*.,leftmargin=1.4em]
  \item \key{Ctrl+A} selects everything on the scene.
  \item Open the \panel{Operations} pane (\menu{Panes>Tools>Operations}). It shows
        one coloured square per operation. Click the \textbf{blue} square to
        assign everything to the \mkseries{Engrave} operation.
  \item Click the rectangle on the scene to select just it, then hold \key{Shift}
        and click the circle. Both are now selected.
  \item Click the \textbf{red} square. The rectangle and the circle move to the
        \mkseries{Cut} operation; the text stays engraved.
  \item Look at the tree. Each operation now lists the elements it contains. If
        something is in the wrong place, drag it from one operation to the other,
        or right-click it and use \menu{Assign Operation}.
\end{enumerate}
\end{walkthrough}

\noindent Text elements belong to \mkseries{Engrave} or \mkseries{Raster}, not
\mkseries{Cut}: a cut operation traces outlines, and letters have no useful
outline to trace into a hole. Everything you need to know about the five
operation types is in Chapter~\ref{ch:ops}.

\begin{warnbox}[Watch for orphans]
If a shape is in the elements branch but referenced by no operation, it will not
burn and MeerK40t will warn you before starting
(``Some shapes are not assigned to any operation\dots''). Always take that
warning seriously --- usually the fix is one click on an operation square.
\end{warnbox}

\section{Step 5 --- Set speed and power}

This is the part of the job that only your machine can answer, but here is a
sane order of operations. Select the \mkseries{Cut} operation and open
\menu{Tools>Properties}.

\begin{center}\small
\begin{tabular}{@{}L{2.9cm}L{3.0cm}L{3.3cm}L{5.1cm}@{}}
\toprule
\textbf{Setting} & \textbf{Typical 40\unit{W} CO\textsubscript{2} on 3\unit{mm}
plywood} & \textbf{Typical diode on 3\unit{mm} plywood} & \textbf{Why} \\
\midrule
\btn{Speed} & 5--10\unit{mm/s} & 3--6\unit{mm/s} & Slower means more energy per
millimetre; below a few mm/s the material chars. \\
\btn{Power} & 60--80\% & 80--100\% & Enough to cut through in one or two passes. \\
\btn{Passes} & 1, then 2 if needed & 2--4 & Two moderate passes cut cleaner than
one aggressive pass. \\
\btn{Kerf compensation} & 0 until you measure it & 0 & Compensate only when you
know your kerf (Chapter~\ref{ch:tuning}). \\
\bottomrule
\end{tabular}
\end{center}

\noindent Then select the \mkseries{Engrave} operation and give it its own
numbers: engraving is a lighter burn at a higher speed, commonly something like
25--35\unit{mm/s} at 20--40\% power for a visible, shallow mark on wood. ``Power''
may be displayed as a 0--1000 scale or as a percentage depending on your device;
the tooltip on the field tells you which.

\begin{tipbox}[Write the numbers down]
Keep a small table on the wall next to the machine: material, thickness, speed,
power, passes, and the date. Two of those lines save you an entire afternoon
three months from now. Chapter~\ref{ch:tuning} shows how to build the same
information systematically with a test grid and the Material Library.
\end{tipbox}

\section{Step 6 --- Position and check before burning}

\begin{walkthrough}{Frame the job, then estimate it}
\begin{enumerate}[label=\arabic*.,leftmargin=1.4em]
  \item In the \panel{Position} pane, make sure the whole tag lies inside the bed
        rectangle on the scene. Anything outside the bed will not be cut, and on
        some machines it will not even be attempted --- or will trigger a
        soft-limit alarm.
  \item On the \panel{Laser-Control} pane, \tabname{Laser} tab, press
        \btn{Outline}. If your machine has a pilot laser or low-power mode, the
        head traces the job's bounding outline so you can see where it will land.
        If it does not, the outline is drawn on screen instead. Either way, this
        is the step that catches the job positioned 40\unit{mm} to the left of
        your material.
  \item Press \btn{Simulate} to open the \panel{Simulation} window, then press
        \btn{Recalculate}. You get an animated preview and time estimates:
        \emph{Laser Time}, \emph{Travel Time}, \emph{Total Time} and distances.
        Watching the animation is the fastest way to see an obviously wrong
        order --- the head leaping back and forth across the bed between two
        cuts that are next to each other.
  \item Close the simulation window when you are satisfied.
\end{enumerate}
\end{walkthrough}

\section{Step 7 --- Run the job}

\begin{walkthrough}{Cut the tag}
\begin{enumerate}[label=\arabic*.,leftmargin=1.4em]
  \item Check the \btn{Optimize} checkbox on the Laser tab. With it on, the job
        is planned through the full pipeline; with it off, the plan skips
        optimisation passes. Beginners should leave it on.
  \item Press \btn{Start}. Deliberately: the button is guarded so that a stray
        keystroke cannot begin a burn, which is why the mouse click matters.
  \item Watch the progress header on the Laser tab. It reports the job state and
        the percentage; the detail line shows steps completed and time remaining.
  \item If the material begins to flare, press \btn{Pause} and be ready with
        \btn{Stop}. \btn{Pause} finishes the current move and holds;
        \btn{Resume} continues. \btn{Stop} aborts the job.
  \item When the job ends, wait for the head to stop and the smoke to clear
        before opening the lid.
\end{enumerate}
\end{walkthrough}

\noindent Behind that one click, the program ran the familiar chain of stages:
copy the operations, preprocess the coordinates, validate them, convert to
cutcode (\emph{blob}), add optimisation commands, optimise, and spool. You can
watch the same stages run line by line in the \panel{Console} pane, and run them
yourself in the \panel{Plan} tab or the \btn{Execute Job} window
(Chapter~\ref{ch:order}). If you ever want the manual version, those are the
places to look.

\begin{behindbox}[Job progress, and what ``Optimize'' really changed]
With \btn{Optimize} ticked, the plan is built and optimised before the spooler
receives it. With it unticked, the same plan is built without the optimisation
stages. The spooler then feeds the driver, and the driver feeds the machine. Each
stage reports into the console channel, which is why the console is the single
best diagnostic tool in the program.
\end{behindbox}

\section{Step 8 --- Judge the result, then save}

Lift the part out and look at the cut edge. Three outcomes:

\begin{itemize}
  \item \textbf{It fell out cleanly.} Save the project (\key{Ctrl+S}) and write
        your settings down. You have a working recipe.
  \item \textbf{The outline cut, the text is faint.} The \mkseries{Engrave}
        operation needs more power or less speed.
  \item \textbf{The edges are charred and the material is sticky.} Too much energy:
        raise the speed or lower the power. On wood, charring often means the
        beam dwelled too long --- one faster pass beats two slow ones.
  \item \textbf{It did not cut through.} Add a pass, then reduce speed, then raise
        power --- in that order, and re-test on scrap each time.
\end{itemize}

\noindent Saving is worth a word, because it is not what most people expect.
MeerK40t's project file is an \file{.svg} file --- the same format you import ---
with extra attributes that preserve operations, settings and notes
(\key{Ctrl+S}, or \menu{File>Save As} for a new name). It is not a proprietary
container, and it is not a device job file; if you want the \file{.egv},
\file{.gcode}, \file{.rd} or \file{.lmc} that the machine itself consumes, use
\menu{Plan>Export} on the Laser panel (Chapter~\ref{ch:files}).

\begin{tryit}{Repeat the exercise on your own material}
Now do it again without reading the steps, with one change: cut a material you
actually care about.
\begin{enumerate}[label=\arabic*.,leftmargin=1.4em]
  \item Draw or import a part with one cut contour and one hole.
  \item Create \mkseries{Cut} and \mkseries{Engrave} operations and assign the
        shapes by colour.
  \item Take a guess at the settings, write the guess down, and cut it.
  \item Adjust exactly one setting and cut again.
  \item Record both attempts. That page of notes is your first material profile,
        and Chapter~\ref{ch:tuning} turns it into something repeatable.
\end{enumerate}
\end{tryit}
